\cotefile{57}
\scanpage{210}
\uline{Proposition} $X$ projectif sur $\unsure{S|Y}$ \inserted{noethérien}, $f_*(\underline{\mathcal O}_X)=\underline{\mathcal O}_Y$ universellement, fibres primaires, $P=\underline{\mathrm{Pic}}_{X/S}$, soit $P'$ la partie de $P$ réunion des $(P_s)'$, où $P'_s$ est \unsure{l'ens des|l'ensemble des} éléments de $P_s^\tau$ image inverse du sous-groupe de $P_s^\tau/P_s^0$ formé des éléments d'ordre premier à \deleted{$p$} $\operatorname{car} k(s)$. \inserted{Supposons que dans $P^\tau$ l'adhérence de \unsure{la $0$|$0$|tt pt} soit quasi-compacte. ou que les $(P_s)^0$ soient \uline{propres sur $k(s)$}.} Alors $P\deleted{_{\unsure{s|S}}}'$ est une partie \deleted{\unsure{fermée, de $P$}} \inserted{\uline{constructible}} \uline{quasi-compacte} de $P$, \deleted{\unsure{ou}}.

\uline{Démonstration}. Il suffit de prouver (réc. noeth sur $Y$) qu'il existe un ouvert non vide $U$ de $Y$, tel que la partie de $P'_{\unsure{s|S}}$ au dessus soit \unsure{qu. cpte.|qu.-cpt} On peut supposer déjà $Y$ affine intègre ; soit $\unsure{g|\eta}$ son pt générique. \deleted{Munissons} $P'_g$ est un sous-groupe \inserted{ouvert} de $P_g$, de type fini sur $k(g)$. Il provient d'un groupe \inserted{$H$} de type fini sur un voisinage de $g$, et on peut supposer \inserted{(rapetissant $Y$)} que ce dernier est un ouvert de\deleted{\unsure{ns}} $P^\tau$, et \unsure{$=$|$\simeq$|à|est} un \uline{ss-groupe ouvert} de $P^\tau$ [prendre un ouvert qu. cpt $V$ de $P^\tau$ contenant $P_g'$].
% archive sheet number: 209
